\documentclass[12pt, titlepage]{article}
\usepackage{amsmath}
\usepackage{amsfonts}
\usepackage{amssymb}
\usepackage{fancyhdr}
\usepackage{cancel}
\usepackage{xcolor}
\usepackage[shortlabels]{enumitem}
\usepackage[bidi=basic-r]{babel}
\babelprovide[main,import,Alph=letters]{hebrew}
\babelfont[hebrew]{rm}{Hadasim CLM}
\babelfont[hebrew]{sf}{Miriam Mono CLM}
\usepackage{titlesec}
\title{קומבינטוריקה 1040286 \\ גליון 3}
\author{יונתן אבידור - 214269565\\רותם עשת - 212674683}
\titleformat{\section}{\Large\sffamily\bfseries}{\thesection}{1em}{}
\newcommand\Ccancel[2][black]{\renewcommand\CancelColor{\color{#1}}\cancel{#2}}
\begin{document}
\maketitle
\pagestyle{fancy}
\fancyhead{}
\fancyfoot[L]{\text{214269565}}
\fancyfoot[R]{\text{212674683}}
\part*{שאלה 1}
\section*{סעיף א'}
מצאו קבועים $a,b$ כך ש-
\[
    k^2 = a\binom{k}{1} + b\binom{k}{2}
\]
\section*{סעיף ב'}
חשבו את הסכום
\[
    \sum_{k=0}^{n}k^2
\]
\newpage
\part*{פתרון 1}
\section*{סעיף א'}
נפתח את הביטוי $a\binom{k}{1} + b\binom{k}{2}$
\begin{gather*}
 a\frac{k!}{1! \left(k-1\right)!} + b\frac{k!}{2!\left(k-2\right)!}\\
    = a \frac{k!}{\left(k-1\right)!} + b\frac{k!}{2!\left(k-2\right)!}\\
    = a \frac{k\cdot\cancel{\left(k-1\right)\cdot\left(k-2\right)\cdots 1}}{\cancel{\left(k-1\right)\cdot\left(k-2\right)\cdots 1}} + b\frac{k\cdot\left(k-1\right)\cdot\cancel{\left(k-2\right)\cdots 1}}{\cancel{\left(k-2\right)\cdots 1} \cdot 2 \cdot 1} \\
    = a \cdot k + b \cdot \frac{k\cdot\left(k-1\right)}{2}
\end{gather*}
נשווה ל-$k^2$
\begin{gather*}
    k^2= a \cdot k + b \cdot \frac{k\cdot\left(k-1\right)}{2}\\
    \implies 2k^2 = 2a \cdot k + b\cdot\left(k^2-k\right) = b\cdot k^2 + \left(2a-b\right) \cdot k
\end{gather*}
קיבלנו שוויון פולינומים, לכן כל מקדם בצד ימין של המשוואה שווה למקדם המתאים בצד שמאל של המשוואה
\[
\begin{cases}
    2 = b \\
    0 = 2a-b
\end{cases}
\]
נציב $b=2$ במשוואה השנייה
\[
0 = 2a -2 \implies 2 = 2a \implies a =1
\]
לכן
\[
    \boxed{a=1,b=2}
\]
\section*{סעיף ב'}
נשים לב שמכיוון ש $0^2 =0$, מתקיים $\sum_{k=0}^n k^2 = \sum_{k=1}^n k^2$\\
נציב את מה שמצאנו בסעיף א'
\begin{gather*}
    \sum_{k=0}^n k^2=   \sum_{k=1}^n k^2 =   \sum_{k=1}^n \binom{k}{1} + 2\binom{k}{2} = \sum_{k=1}^n \binom{k}{1} + 2\sum_{k=1}^n \binom{k}{2}
\end{gather*}
אין דרכים לבחור שני איברים מקבוצה בגודל $1$ ולכן $\binom{1}{2}=0$, מכאן ש
    \[
\sum_{k=1}^n \binom{k}{2}=0+\sum_{k=2}^n \binom{k}{2} = \sum_{k=2}^n\binom{k}{2}
\]
נזכר ב"נוסחת מקל ההוקי" שראינו בתרגול
\[
\sum_{i=r}^n \binom{i}{r} = \binom{n+1}{r+1}
\]
נשתמש בה על שני המחוברים
\begin{gather*}
    \sum_{k=1}^n \binom{k}{1} + 2\sum_{k=2}^n \binom{k}{2} = \binom{n+1}{2}+ 2\binom{n+1}{3}
\end{gather*}
נפתח את נוסחאות הבינום ונפשט
\begin{gather*}
    \binom{n+1}{2} + 2\binom{n+1}{3} = \frac{\left(n+1\right)!}{2\left(n-1\right)!} + \cancel{2}\frac{\left(n+1\right)!}{3\cdot\cancel{2}\left(n-2\right)!} \\
    =    \frac{\left(n+1\right)!}{2\left(n-1\right)!} + \frac{\left(n+1\right)!}{3\left(n-2\right)!} = \frac{n\left(n+1\right)}{2} + \frac{\left(n-1\right)\cdot n\cdot \left(n+1\right)}{3}\\
    \frac{3\left(n\left(n+1\right)\right)+2\left(n\left(n+1\right)\right)\left(n-1\right)}{6} = \frac{\left(n\left(n+1\right)\right)\left(3+2\left(n-1\right)\right)}{6}\\
    = \boxed{\frac{n\left(n+1\right)\left(2n+1\right)}{6}}
\end{gather*}
\newpage
\part*{שאלה 2}
השתמשו בנוסחת המולטינום על מנת למצוא את המספרים הבאים.
\section*{סעיף א'}
מצאו את המקדם של $x^6y^8z^3$ בפיתוח של $\left(x^2+y^2+z\right)^{10}$
\section*{סעיף ב'}
כמה מחוברים יש בפיתוח המולטינומי של $\left(x+y+z\right)^{30}$
\part*{פתרון 2}
\section*{סעיף א'}
נוסחת המולטינום היא
\[
    \left(x_1+x_2+\ldots+x_m\right)^n = \sum_{k_1+k_2+\ldots+k_m=n} \binom{n}{k_1,k_2,\ldots,k_m}\cdot x_1^{k_1} \cdots x_m^{k_n}
\]
לכן
\[
    \left(x^2+y^2+z\right)^{10} = \sum_{k_1+k_2+k_3=10} \binom{10}{k_1,k_2,k_3} \cdot x^{2k_1} \cdot y^{2k_2} \cdot z^{k_3}
\]
אנחנו מחפשים את המקדם של $x^6y^8z^3$\\
נתחיל בלמצוא את $k_{1-3}$
\begin{gather*}
    x^{2k_1} = x^6 \implies 2k_1 = 6 \implies k_1 = 3\\
    y^{2k_2} = y^8 \implies 2k_2 = 8 \implies k_2 = 4 \\
    z^{k_3} = z^3 \implies k_3 = 3
\end{gather*}
לכן המקדם הוא $\binom{10}{3,4,3}$
\[
    \binom{10}{3,4,3} = \frac{10!}{3!\cdot4!\cdot 3!} = \frac{10\cdot\Ccancel[teal]{9}\cdot\Ccancel[orange]{8}\cdot7\cdot\Ccancel[green]{6}\cdot5\cdot\Ccancel[blue]{4\cdot3\cdot2}}{\Ccancel[green]{3\cdot2}\cdot \Ccancel[blue]{4\cdot3\cdot2}\cdot\Ccancel[teal]{3}\cdot\Ccancel[orange]{2}} = 10 \cdot 3 \cdot 4 \cdot 7 \cdot 5 = \boxed{4200}
\]
\section*{סעיף ב'}
לפי נוסחת המולטינום, הפיתוח של $\left(x+y+z\right)^{30}$ הוא
\[
    \sum_{k_1+k_2+k_3=30} \binom{30}{k_1,k_2,k_3}x^{k_1}\cdot y^{k_2} \cdot z^{k_3}
\]
לפי הגדרה, הסכום רץ על כל הצירופים $k_1,k_2,k_3$, כך ש-$k_1+k_2+k_3=30$, כשעבור כל צירוף, יש מחובר אחד, מכאן  שכמות המחוברים שווה לכמות הפתרונות למשוואה $k_1+k_2+k_3=30$ כאשר $k_1,k_2,k_3$ הם שלמים אי-שליליים\\
על פי נוסחה שלמדנו, כמות הפתרונות למשוואה כזו היא $\binom{30+3-1}{3-1}$
\[
    \binom{30+3-1}{3-1}= \binom{32}{2}  =\frac{32!}{2! \cdot 30!} = \frac{32\cdot 31 \cdot \cancel{30 \cdots 2}}{2 \cdot \cancel{30 \cdots 2}}= \frac{32\cdot31}{2} = \frac{992}{2} = \boxed{496} 
\]
\newpage
\part*{שאלה 3}
יש כיתה של $3n$ תלמידים. בכמה דרכים אפשר לחלק אותם לשלשות (לא מסומנות)?\\
השתמשו בנוסחת סטירלינג
\part*{פתרון 3}
על מנת לחלק את הקבוצה לשלשות, נבחר $3$ תלמידים עבור השלשה הראשונה, ואז עוד $3$ עבור השלשה השלישית וכן הלאה $n$ פעמים, עד שכולם מחולקים לשלשות
\begin{gather*}
    \binom{3n}{3} \cdot \binom{3n-3}{3} \cdots \binom{6}{3}\cdot \binom{3}{3}\\
    \frac{\left(3n\right)!}{3!\cdot\cancel{\left(3n-3\right)!}}\cdot \frac{\cancel{\left(3n-3\right)!}}{3!\cancel{\left(3n-6\right)!}}\cdots \frac{\cancel{6!}}{3!\cdot\cancel{3!}}\cdot\frac{\cancel{3!}}{3!} \\
    = \frac{\left(3n\right)!}{\left(3!\right)^n}
\end{gather*}
כעת, על מנת "להיפטר" מהסדר בין השלשות, נחלק בכמות הדרכים לסדר את השלשות השונות ($n!$)
\[
    \frac{\frac{\left(3n\right)!}{\left(3!\right)^n}}{n!}= \frac{\left(3n\right)!}{6^n\cdot n!}
\]
נשתמש בנוסחת סטירלינג על מנת למצוא קירוב לביטוי\\
נתחיל עם המונה
\[
    \left(3n\right)! \approx \left(\frac{3n}{e}\right)^{3n} \sqrt{6\pi n} 
\]
נעבור למכנה
\[
    6^n \cdot n! \approx  6^n \cdot \left(\frac{n}{e}\right)^n \sqrt{2\pi n} 
\]
וביחד
\begin{gather*}
    \frac{\left(3n\right)!}{6^n\cdot n!} \approx  \frac{\left(\frac{3n}{e}\right)^{3n} \sqrt{6\pi n}}{6^n \cdot \left(\frac{n}{e}\right)^n \sqrt{2\pi n} }\\ \\
= \frac{27^{n}\cdot n^{3n} \cdot \frac{1}{e^{3n}} \cdot \left(\sqrt{3}\cdot\sqrt{2\pi n}\right)}{6^n \cdot n^n \cdot \frac{1}{e^n} \cdot \sqrt{2\pi n}} \\ \\
=    \frac{27^n}{6^n} \cdot \frac{n^{3n}}{n^n} \cdot \frac{e^{n}}{e^{3n}} \cdot \frac{\sqrt{3}\cdot \sqrt{2\pi n}}{\sqrt{2\pi n}} \\ \\
= \left(\frac{9}{2}\right)^n \cdot n^{2n} \cdot \frac{1}{e^{2n}} \cdot \sqrt{3} \\\\
= \sqrt{3} \cdot \left(\frac{9n^2}{2e^2}\right)^{n}
\end{gather*}
\newpage
\part*{שאלה 4}
העריכו בעזרת סטרלינג את העצרת הכפולה $\left(2n\right)!!$.
\part*{פתרון 4}
לפי הגדרת עצרת כפולה
\begin{gather*}
  \left(2n\right)!! =2n\cdot \left(2n-2\right) \cdot \left(2n-4\right) \cdots 2 \\
= \left(2\cdot n\right) \cdot \left(2\left(n-1\right)\right) \cdot \left(2\left(n-2\right)\right) \cdots \left(2 \cdot 1\right)\\
=  \underbrace{2\cdot 2 \cdots 2}_{n} \cdot \left(n-1\right)\cdot\left(n-2\right)\cdots1 \\
= 2^n \cdot n!
\end{gather*}
נשתמש בנוסחת סטירלינג
\[
    2^n \cdot n! \approx 2^n \cdot \left(\frac{n}{e}\right)^n \cdot \sqrt{2\pi n} =\boxed{\left(\frac{2n}{e}\right)^n \cdot \sqrt{2 \pi n}}
\]
\newpage
\part*{שאלה 5}
\section*{סעיף א'}
מהו סכום כל המקדמים המולטינומים $\binom{30}{a,b,c}$
\section*{סעיף ב'}
ציינו איזה מבין המקדמים האלה הוא הכי גדול
\section*{סעיף ג'}
מה יותר גדול, סכום המקדמים כנ"ל עם $b$ זוגי, או סכום אלו עם $b$ אי זוגי? מה ההפרש?
\section*{סעיף ד'}
חשבו את גודל הקבוצה
\[
    \left\{\left(x_1,\ldots,x_n\right)\in\left\{0,1,2\right\}^{n}|\text{אפס מופיע מספר זוגי של פעמים }\right\}
\]
\part*{פתרון 5}
\section*{סעיף א'}
נחפש את סכום כל המקדמים $\binom{30}{a,b,c}$
\[
    \sum_{a+b+c=30} \binom{30}{a,b,c} 
\]
נכפיל את כל הביטוי ב-$1^a\cdot1^b\cdot1^c$, נשים לב כי $1^a\cdot1^b\cdot1^c=1\cdot1\cdot1=1$, לכן הכפלה שכזו לא משנה את ערך הביטוי
\[
    \sum_{a+b+c=30} \binom{30}{a,b,c}=   1^a \cdot 1^b \cdot 1^c
    \sum_{a+b+c=30} \binom{30}{a,b,c} 
\]
לפי נוסחת המולטינום, זה שווה לפתרון של המשוואה $\left(1+1+1\right)^{30}$, לכן
\[
\sum_{a+b+c=30} \binom{30}{a,b,c}  = \left(1+1+1\right)^{30} = \boxed{3^{30}}
\]
\section*{סעיף ב'}
המקדם הגדול ביותר יהיה המקדם בעל המכנה הקטן ביותר, שכן לכל מקדם המונה הוא $30!$\\
לכן הבעיה שקולה ללמצוא בחירה ל-$a,b,c$ כך ש-$a+b+c=30$ שמביאה את הערך המינמלי לביטוי $a!b!c!$\\
נטען כי הבחירה הזו היא $a=b=c=10$
\subsection*{למה}
יהיו $x\leq y$ אזי
\[
    x!y! \leq \left(x-1\right)!\left(y+1\right)!
\]
\subsection*{הוכחת הלמה}
נסתכל על היחס ביניהם, ונרצה להראות כי הוא קטן מ-$1$
\begin{gather*}
    \frac{x!y!}{\left(x-1\right)!\left(y+1\right)!} = \frac{\Ccancel[red]{1\cdot2\cdots\left(x-1\right)}\cdot x\cdot \Ccancel[green]{1\cdot2\cdots y}}{\Ccancel[red]{1\cdot2\cdot\left(x-1\right)}\cdot\Ccancel[green]{1\cdot2\cdots y}\cdot\left(y+1\right)} \\
    = \frac{x}{y+1} \leq \frac{y}{y+1} < 1
\end{gather*}
כנדרש\\
$\blacksquare$
\subsubsection*{בחזרה להוכחה}
נשים לב שכל צירוף של $a,b,c$ כך ש-$a+b+c=30$ מתקבל על ידי לקיחת הסידור $a=b=c=10$ ואז הורדת $1$ מאחד האיברים והוספת $1$ לאיבר אחר הגדול או שווה לו בגודלו כמות פעמים כלשהי.\\
זה מכיוון שעל מנת שהסכום הכולל ישאר $30$, על כל $1$ שמתווסף לאחד מהאיברים, יש להוריד $1$ מאחד אחר.\\
נראה שאפשר לעשות זאת בלי להעביר מאיבר אחד לאיבר הקטן ממש ממנו\\
נסמן בצירוף כלשהו של $a,b,c$ ב-$x$ את האיבר הקטן ביותר, $y$ האיבר השני הכי קטן, ו-$z$ האיבר הגדול ביותר\\
על מנת להפוך את $10$ ל-$x$, נחסיר כל פעם $1$ ונוסיף פעם אחת ל-$y$ ופעם אחת ל-$z$, כאשר אם יש לעשות כמות אי-זוגית של העברות, את האחת הנוספת נעביר ל-$z$. לאחר מכן נעשה את כמות ההעברות מ-$y$ ל-$z$, ובכך נמנענו מלעשות העברה מאיבר מסויים לאיבר קטן ממש ממנו.\\
הראינו כי כל העברה שנעשה מ-$a=b=c=10$ תגדיל את $a!b!c!$ ובכך תקטין את הביטוי $\frac{30!}{a!b!c!}$\\
מכאן שהצירוף הגדול ביותר הוא 
\[
    \boxed{a=10,b=10,c=10}
\]
\section*{סעיף ג'}
סכום המקדמים עם $b$ זוגי מתאר את כמות הדרכים לחלק $30$ איברים סדורים לקבוצות $a,b,c$ כך ש-$a+b+c=30$ וכן $b$ זוגי. הדבר שקול למספר הסדרות באורך $30$ מעל $a,b,c$ כך שמספר המופעים של $b$ הינו זוגי\\
נטען כי כמות הסדרות בהן $b$ מופיע כמות זוגית של פעמים גדולה ב-$1$ מכמות הסדרות בהן $b$ מופיע כמות אי-זוגית של פעמים\\
נקרא לסדרה בעלת כמות זוגית של $b$ "סדרה זוגית", ולסדרה בעלת כמות אי-זוגית של $b$ "סדרה אי-זוגית" \\
נראה באינדוקציה על $n$ כי כמות הסדרות הזוגיות באורך $n$ גדולה ב-$1$ מכמות הסדרות האי-זוגיות באורך $n$
\subsection*{בסיס האינדוקציה}
עבור $n=1$ מתקבלות הסדרות הבאות:\\
$\left(a\right)$ (סדרה זוגית) \\
$\left(b\right)$ (סדרה אי-זוגית) \\
$\left(c\right)$ (סדרה זוגית) \\
קיבלנו שיש $2$ סדרות זוגיות וסדרה זוגית אחת
\subsection*{הנחת האינדוקציה}
נניח שעבור סדרה באורך $n$ קיימות $k$ סדרות אי-זוגיות ו-$k+1$ סדרות זוגיות.
\subsection*{צעד האינדוקציה}
כל סדרה באורך $n+1$ מתקבלת על ידי הוספה של איבר כלשהו מבין $\left\{a,b,c\right\}$ לסוף של סדרה קיימת באורך $n$\\
עבור הוספה של האיבר $b$ לסוף סדרה באורך $n$, מתקבלות $k$ סדרות זוגיות ו-$k+1$ סדרות אי-זוגיות (מהנחת האינדוקציה והעובדה שהוספת $b$ משנה את הזוגיות של סדרה)\\
עבור הוספה של האיברים $a$ או $c$ לסוף סדרה באורך $n$, מתקבלות $k+1$ סדרות זוגיות ו-$k$ סדרות אי-זוגיות (מהנחת האינדוקציה והעובדה שלהוסיף $a$ או $c$ לא משנה את הזוגיות של הסדרה). אלו כל הדרכים לבנות סדרה באורך $n+1$ מעל $\left\{a,b,c\right\}$\\
נסכום את כמות הסדרות הזוגיות באורך $n+1$
\[
    \left(k\right) + \left(k+1\right) + \left(k+1\right) = 3k+2
\]
נסכום את כמות הסדרות האי-זוגיות באורך $n+1$
\[
    \left(k+1\right) + \left(k\right) + \left(k\right) = 3k+1
\]
ואכן קיבלנו כי כמות הסדרות הזוגיות גדולה ב-$1$ מכמות הסדרות האי-זוגיות\\
לכן סכום המקדמים כאשר $b$ זוגי גדול מהסכום כאשר $b$ אי-זוגי, באחד.
\section*{סעיף ד'}
הראינו בסעיף הקודם כי כמות הסדרות הזוגיות (במקרה הזה נגדיר סדרה זוגית לפי כמות המופעים של $0$) גדולה ב-$1$ מכמות הסדרות האי-זוגיות הללו.\\
מכיוון שאלו שני המצבים היחידים, הסכום שלהם הוא מספר האפשרויות עבור סדרות כאלה ללא תנאי על זוגיות הסדרות\\
כעת, מכיוון שכמות הסדרות באורך $n$ הוא $3^n$ (מספר אי זוגי). מתקבל כי כמות הסדרות הזוגיות באורך $n$ הוא
\[
    \boxed{  \lfloor \frac{3^n}{2}\rfloor+1 = \frac{3^n+1}{2}}
\]
\newpage
\part*{בדיחת קרש}
איך קוראים לחתול שעושה שיעורי בית ביום ההולדת שלו\\
מיציתי
\end{document}
